\section{Assunzione impegni}
Con la sottoscrizione del presente documento, ciascun componente del gruppo \textit{\textbf{Git Happens}} si impegna formalmente a dedicare allo svolgimento del progetto un monte orario complessivo di 93 ore produttive ciascuno.
Per garantire una formazione omogenea e un contributo equamente bilanciato, tutti i ruoli tipici del ciclo di vita del software saranno coperti a rotazione da ogni membro del team per una durata equivalente.
I ruoli e le relative responsabilità sono così definiti:

\begin{itemize}
    \item \textbf{Responsabile}: Figura di riferimento per l'intero team; si occupa della pianificazione delle attività, della gestione delle risorse, del coordinamento del lavoro e del rispetto dei tempi e dei vincoli stabiliti.
    \item \textbf{Amministratore}: Gestisce i processi, le configurazioni e le risorse operative. Monitora l'avanzamento delle attività, gestisce i cambiamenti e garantisce l'efficienza dei processi di sviluppo.
    \item \textbf{Analista}: Si occupa dell'analisi dei requisiti, raccogliendo le esigenze degli \textit{stakeholder} e traducendole in specifiche dettagliate per il team di sviluppo.
    \item \textbf{Progettista}:  Definisce l'architettura complessiva del sistema, identificandone moduli, componenti e interazioni a partire dalle specifiche definite dagli analisti.
    \item \textbf{Programmatore}: Responsabile dell'implementazione e della scrittura del codice sorgente in conformità alle specifiche di progetto.
    \item \textbf{Verificatore}: Cura il controllo qualità di codice e documentazione, verificando il rispetto degli standard definiti, la conformità ai requisiti e l'assenza di difetti.
\end{itemize}